\begin{theorem}[\statusfalse: Theorem 2.1 of \cite{McSwiggenSahi2026}]\label{thm:macdonald}
The stated Schur-convexity fails already in rank two.
\end{theorem}
\begin{proof}
Set $q=t=r\in(0,1)$, $a=1$, and choose lattice index $\nu=(1,0)$.  The resulting lattice point is $x=(1,1)$.  Let $\lambda=(2,0)$ and $\mu=(1,1)$; then $\lambda\succeq\mu$.  At $q=t$, Macdonald polynomials specialize to Schur polynomials, and direct evaluation gives
\[
\Omega_{(2,0)}(1,1;r,r)=\frac{3}{1+r+r^2},\qquad
\Omega_{(1,1)}(1,1;r,r)=\frac1r.
\]
But
\[
1+r+r^2-3r=(1-r)^2>0,
\]
so $\Omega_{(2,0)}<\Omega_{(1,1)}$ for every $0<r<1$, opposite to Schur-convexity.  At $r=1/2$, the values are $12/7<2$.
\end{proof}
